% ============================================================
% BME VIK TDK abstract — VARIANT 4: the combined stack (F2 + F3 + F4)
%
% Fourth of the framings. This is the thesis-scale version: it keeps
% variant 3's objective (novelty as a generative residual) but puts
% the two results it depends on in front of it, in the order in which
% they have to be established.
%
%   Stage 1 — the model, and its calibration on the part of a paper
%     the bibliography DOES determine, measured without an LLM judge.
%   Stage 2 (was variant 2) — attribution. Is that recovery a fact
%     about graph structure, or only about having more text?
%   Stage 3 (was variant 3) — the part the bibliography does not
%     determine, read as novelty.
%
% ORDER: model first, then attribution, then novelty. An earlier draft
% led with the controls; that forward-referenced the prefix before the
% architecture was introduced, and made the controls read as a
% defensive gate rather than as attribution of a result already shown.
%
% The load-bearing claim is that these are not three projects: stages
% 2 and 3 are the two halves of one residual, and stage 1 is what
% makes either of them a statement about the graph.
%
% Structure follows the official one-page sample:
% https://portal.vik.bme.hu/files/00000047.pdf
%   - max 1 page, min ~20 lines of body text
%   - NO header, footer, or page numbers
%   - figures/equations only if strictly necessary
%   - own contribution must be clearly separated from cited work
%   - body must cover: methodology, key results, conclusions
%
% NUMBERS: three slots, one per stage — the margin between the
% graph-conditioned model and the text-matched baseline; the share of
% downstream performance recovered from generated node text; and the
% validation figures for the novelty measure. Do not submit without
% them; the official sample requires key results in the body.
%
% LENGTH WARNING: this variant carries three stages in the space v1
% and v2 use for one. It is at the page limit. If it overflows after
% translation to Hungarian (which runs longer), cut from stage 1 —
% the control suite can be named without being enumerated.
%
% NOTE ON LANGUAGE: the official TDK conference language is Hungarian
% (English talks require prior agreement with the section chair, per
% the official sample). This English draft is for internal review with
% Bálint/Dani; translate to Hungarian before submission unless an
% English exception is confirmed.
% ============================================================

\documentclass[12pt]{article}

\usepackage[utf8]{inputenc}
\usepackage[T1]{fontenc}
\usepackage[english]{babel}
\usepackage[a4paper, margin=2.5cm]{geometry}
\usepackage{mathptmx}      % Times-like font
\usepackage{hyperref}
\usepackage{enumitem}

\pagestyle{empty}          % <- enforces "no header/footer/page numbers"
\setlength{\parindent}{0pt}
\setlength{\parskip}{0.6em}

\begin{document}

% ---------- Title ----------
\begin{center}
    {\Large\bfseries What a Citation Graph Determines, and What It Does Not: \\
    From Graph-Conditioned Language Models to a Measure of Novelty}
\end{center}

\vspace{0.8em}

% ---------- Authors ----------
\begin{center}
    Gema Csanád IV. Inf, gemacsanad@gmail.com \\

\end{center}

\vspace{0.4em}

% ---------- Consultant ----------
\begin{center}
    Consultant: Unyi Dániel
\end{center}

\vspace{1.2em}

% ==========================================================
% BODY
% ==========================================================

A citation graph represents scientific literature as a network: each
node is a publication, each edge a citation from one paper to
another, and every node carries that paper's abstract as its text
content. A reference list records the problem a paper attacks, the
methods it borrows and the community it belongs to. We ask how much
of a paper that record determines — and we use both halves of the
answer: what a citation neighbourhood predicts about a paper is
recoverable data, and what it fails to predict is a measure of the
paper's novelty.

Methodologically, a graph neural network encodes the citation
neighbourhood — the cited papers, and their references in turn,
beyond what could fit in a prompt — into a continuous representation,
which a trained adapter projects into soft prefix tokens injected
into the attention layers of a language model. The model generates
the target paper's abstract from this representation together with
the plain text of the directly cited papers, never seeing the
target's own words. The target's features are masked and edges run
from cited to citing paper, so nothing published after the target can
reach it.

How much of a paper this recovers is measured without a language
model as judge. The generated text is substituted for the real node
attributes, and we measure how much of the downstream performance
that real text supports — node classification on the citation graph —
survives the substitution. This is an objective criterion, and it
answers a practical question in its own right: which nodes in a graph
with missing text can be restored from their neighbourhood, and which
can only be described in outline.

Whether that recovery is a fact about graph structure, or merely
about having more text, is a separate question, and the one existing
work leaves open. Structure is usually given to a language model by
encoding a neighbourhood into a continuous vector and injecting it
alongside the neighbours' text, and gains are then credited to the
structure. But text is the cheaper channel, and a model offered the
same information twice has no reason to learn the harder one; prior
graph-conditioned models report end-task performance only, so this is
never tested. We test it against a text-matched baseline receiving
the same neighbourhood content as plain text within the same token
budget, and against prefixes built from a disabled or a shuffled
neighbourhood, which separate the prefix's content from its presence.

The remainder is the measurement. A paper's contribution is by
definition its departure from the work it cites, so the discrepancy
between the predicted and the actual abstract estimates its novelty.
Reference-based measures score the atypicality of cited journals and
discard content; perplexity-based measures ask whether a text is
unusual in general rather than given this paper's own antecedents.
Here the expectation is built from the paper's own reference list.
Because such a residual confounds novelty with model error and with
incomplete coverage, it is read relative to matched papers of the
same field, year and reference count — and only against a model whose
accuracy where prediction is possible has already been established.

Our own contribution is this argument and the system that carries it:
the judge-free calibration of what a citation neighbourhood recovers,
the controls that decide whether that recovery is structural, and the
formulation and validation of novelty as a conditional generative
residual. Novelty has no ground truth, so validation is the
substantive difficulty; we approach it through axiomatic properties
any such measure should satisfy, separation of recognised landmark
papers from matched controls, and correlation with the disruption
index and later citation impact on the same graph. One quantity, two
uses: what a bibliography predicts is data that can be restored, and
what it fails to predict is what makes a paper worth reading.

\end{document}
